\begin{lemma}
Let $g\in\mathsf{IncSeq}$ and let $k\in\mathbb N$ satisfy
$k<g(k)$ and $g(j)=j$ for every $j<k$.  For every
$f\in\mathsf{IncSeq}$ and every $l\in\mathbb N$, define
$\sigma_f(l)=\mathsf{BlockSigma}(g,k,f,l)$ as in
\noderef{BlockSigma}.  The block maps satisfy
\[
\sigma_{f\circ\mathsf{SuccSeq}}(l)=\sigma_f(g(l)).
\]
Here $f\circ\mathsf{SuccSeq}$ is the right composition represented by
$\mathsf{RightComp}(\mathsf{SuccSeq},f)$, and equivalently
$\sigma_{f\circ\mathsf{SuccSeq}}=\sigma_f\circ g$.
\end{lemma}
\begin{proof}
Put $G(n)=\mathsf{OrbitPoint}(g,k,n)=g^{[n]}(k)$, using
\noderef{OrbitPoint}; in particular $G(0)=k$.  By
\noderef{OrbitBlockCoverage}, every $l$ either satisfies $l<G(0)$ or
belongs to a unique block $[G(n),G(n+1))$.

First suppose that $l<G(0)$.  Then $l<k$, so the hypothesis that $g$ fixes
every point below $k$ gives $g(l)=l$.  The initial branch of the definition
in \noderef{BlockSigma} therefore gives
\[
\mathsf{BlockSigma}(g,k,\mathsf{RightComp}(\mathsf{SuccSeq},f),l)=l.
\]
Also $g(l)=l<G(0)$, so the same initial branch gives
\[
\mathsf{BlockSigma}(g,k,f,g(l))=l.
\]
Thus the two sides are equal in this case.

Now suppose that $l$ is not below $G(0)$.  Choose the unique $n$ supplied
by \noderef{OrbitBlockCoverage}, so that
\[
G(n)\leq l<G(n+1).
\]
Since $g\in\mathsf{IncSeq}$, \noderef{IncSeq} makes $g$ an order
embedding and hence strictly increasing.  Applying this first to the left
inequality and then to the right inequality yields
\[
g(G(n))\leq g(l)<g(G(n+1)).
\]
We next identify both endpoints.  From the iterate definition in
\noderef{OrbitPoint}, the successor rule for function iterates gives
\[
G(n+1)=g^{[n+1]}(k)=g(g^{[n]}(k))=g(G(n)).
\]
Applying the same calculation with $n+1$ in place of $n$ gives
\[
G(n+2)=g^{[n+2]}(k)=g(g^{[n+1]}(k))=g(G(n+1)).
\]
Consequently
\[
G(n+1)\leq g(l)<G(n+2),
\]
so $g(l)$ lies in the $(n+1)$st orbit block.  This explicitly locates the
argument of the right-hand block map before applying its defining formula.

By \noderef{RightComp} and \noderef{SuccSeq},
\mathsf{RightComp}(\mathsf{SuccSeq},f)(n)=f(n+1)$.  Therefore the block
branch of \noderef{BlockSigma}, applied to the original block of $l$, gives
\[
\mathsf{BlockSigma}(g,k,\mathsf{RightComp}(\mathsf{SuccSeq},f),l)
=g^{[f(n+1)-n]}(l).
\]
The block branch applied to the $(n+1)$st block of $g(l)$ gives
\[
\mathsf{BlockSigma}(g,k,f,g(l))
=g^{[f(n+1)-(n+1)]}(g(l)).
\]

It remains to compare these iterates.  We first prove the needed index
bound for every increasing embedding $f$.  For $m=0$, $0\leq f(0)$ because
$f(0)$ is a natural number.  If $m\leq f(m)$, then strict increase from
\noderef{IncSeq} gives $f(m)<f(m+1)$, hence
$m+1\leq f(m)+1\leq f(m+1)$.  Induction proves
\[
\forall m\in\mathbb N,\qquad m\leq f(m).
\]
In particular, $n+1\leq f(n+1)$.  Set
\[
a=f(n+1)-(n+1).
\]
The preceding inequality makes this subtraction exact, and elementary
natural-number arithmetic gives
\[
a+1=(f(n+1)-(n+1))+1=f(n+1)-n.
\]

We also prove the iterate shift used here rather than treating it as
implicit.  For every $a$ and $x$, function-iterate recursion gives
\[
g^{[a]}(g(x))=g^{[a+1]}(x).
\]
Indeed, for $a=0$ both sides are $g(x)$.  If the equality holds for $a$,
then applying $g$ to both sides and using the successor rule for iterates
gives
\[
g^{[a+1]}(g(x))
=g(g^{[a]}(g(x)))
=g(g^{[a+1]}(x))
=g^{[a+2]}(x),
\]
which is the claim for $a+1$.  Taking $x=l$ and the above value of $a$ now
gives
\[
g^{[f(n+1)-(n+1)]}(g(l))
=g^{[a]}(g(l))
=g^{[a+1]}(l)
=g^{[f(n+1)-n]}(l).
\]
This is exactly the equality of the two block values, and the two cases
prove the lemma.
\end{proof}
